% chapters/17/section-04.tex
\section{Lembar Refleksi Akhir Semester}
\label{sec:refleksi-akhir}

\subsection{Tiga Pertanyaan yang Harus Dijawab}

Lembar refleksi akhir semester memuat tiga pertanyaan. Ketiganya tampak
sederhana, dan ketiganya sering dijawab dengan kalimat umum yang tidak
mengandung apa pun. Bagian ini menjelaskan apa yang dicari dari setiap
jawaban, agar Anda dapat menjawabnya dengan berguna.

\begin{enumerate}
  \item \textbf{Konsep apa yang paling sulit dipahami pada blok Python/Pygame
        dan blok Unity, dan bagaimana usaha yang dilakukan untuk
        mengatasinya?}
        \emph{Yang dicari:} konsep yang spesifik, bukan ``Unity itu sulit''.
        Sebutkan konsepnya, sebutkan gejala salah pahamnya, dan sebutkan
        tindakan yang Anda lakukan beserta hasilnya.
  \item \textbf{Apa perbedaan pengalaman bekerja dengan Python/Pygame
        dibandingkan dengan Unity, ditinjau dari kemudahan, kendali atas
        algoritma, dan kecepatan pengembangan?}
        \emph{Yang dicari:} jawaban yang memakai ketiga tinjauan itu secara
        terpisah, disertai contoh dari pekerjaan Anda sendiri. Sebuah jawaban
        yang menyatakan salah satu lebih baik tanpa memisahkan tinjauannya
        belum menjawab pertanyaan ini.
  \item \textbf{Kompetensi apa yang ingin dikembangkan setelah mata kuliah
        ini, dan bagaimana rencana untuk mencapainya?}
        \emph{Yang dicari:} satu atau dua kompetensi yang nyata, beserta
        langkah pertama yang dapat Anda mulai dalam sebulan ke depan.
\end{enumerate}

\subsection{Menjawab dengan Bukti, Bukan dengan Kesan}

Refleksi yang dinilai tinggi memiliki satu ciri yang sama: setiap pernyataannya
dapat ditelusuri ke sesuatu yang benar-benar terjadi. Perbandingan keduanya
diperlihatkan pada Tabel~\ref{tab:refleksi-bukti}.

\begin{table}[htbp]
\centering
\caption{Jawaban berdasarkan kesan dan jawaban berdasarkan bukti}
\label{tab:refleksi-bukti}
{\setlength{\tabcolsep}{4pt}\footnotesize%
\begin{tabularx}{\textwidth}{@{}L{4.2cm} Y Y@{}}
\toprule
\headrow \textbf{Pertanyaan} & \textbf{Jawaban berdasarkan kesan} & \textbf{Jawaban berdasarkan bukti} \\
\midrule
Konsep tersulit &
``\istilah{Delta time} agak membingungkan.'' &
``Pada Bab~IV, pemain saya bergerak lebih cepat ketika jendelanya diperbesar,
sebab saya memakai jumlah bingkai, bukan \istilah{delta time}. Saya
memperbaikinya dengan mengalikan kecepatan dengan selisih waktu, dan
mengujinya pada dua ukuran jendela.'' \\
\addlinespace
Perbedaan Pygame dan Unity &
``Unity lebih mudah.'' &
``Di Pygame saya menulis sendiri pemeriksaan tabrakan per sumbu, sehingga saya
tahu persis mengapa pemain berhenti di dinding. Di Unity saya tidak menulisnya,
dan ketika pemain tersangkut di sudut, saya perlu membuka panel
\istilah{Physics Debug} untuk memahami sebabnya.'' \\
\addlinespace
Rencana lanjutan &
``Ingin belajar lebih banyak tentang game.'' &
``Saya ingin memahami animasi yang digerakkan kerangka, sebab karakter saya
sekarang hanya memakai potongan gambar. Langkah pertama: mengerjakan satu
tutorial karakter berkerangka pada bulan depan, lalu membandingkannya dengan
animasi potongan gambar yang sudah saya buat.'' \\
\bottomrule
\end{tabularx}}
\end{table}

Jawaban pada kolom ketiga selalu memuat tiga hal: keadaan
yang konkret, tindakan yang diambil, dan hasil atau rencana yang dapat
diperiksa. Ketiga hal itu pula yang membuat refleksi berguna bagi Anda
sendiri, bukan hanya bagi penilaian.

\subsection{Lembar yang Dipakai}

Bentuk lembar refleksi yang dipakai pada bab ini diperlihatkan pada
Tabel~\ref{tab:lembar-refleksi}. Isilah dengan tangan atau pada berkas, dan
serahkan bersama kelengkapan Proyek Akhir.

\begin{table}[htbp]
\centering
\caption{Lembar refleksi akhir semester}
\label{tab:lembar-refleksi}
{\setlength{\tabcolsep}{4pt}\footnotesize%
\begin{tabularx}{\textwidth}{@{}L{4.0cm} Y@{}}
\toprule
\headrow \textbf{Bagian} & \textbf{Isi} \\
\midrule
Konsep tersulit di blok Python & \\[1.6em]
\midrule
Konsep tersulit di blok Unity & \\[1.6em]
\midrule
Usaha yang dilakukan, dan hasilnya & \\[1.6em]
\midrule
Perbedaan ditinjau dari kemudahan & \\[1.4em]
\midrule
Perbedaan ditinjau dari kendali algoritma & \\[1.4em]
\midrule
Perbedaan ditinjau dari kecepatan pengembangan & \\[1.4em]
\midrule
Kompetensi yang ingin dikembangkan & \\[1.4em]
\midrule
Langkah pertama dan waktunya & \\[1.8em]
\bottomrule
\end{tabularx}}
\end{table}
